\subsubsection*{2.4.5.}
设 $x\in X_s$ 是 $(X,F)$ 不光滑的一个孤立点. 更一般地, 设 $K\in\operatorname{Ob}D^b(X,\Lambda)$, $x$ 为 $R\Phi(K)$ 支集的孤立点. 设 $i:\{x\}\hookrightarrow X_s$ 为包含映射. 复形 $i_*i^*R\Phi(K)$ 在 $x$ 的一个邻域中同构于 $R\Phi(K)$, 且由 2.4.3, 它是 $R\Phi(K)$ 的一个直和项. 因此, 对 $\sigma\in I$, 变差态射诱导
\[
\operatorname{Var}(\sigma):
i_*i^*R\Phi_{\bar\eta}(K)\longrightarrow R\Psi_{\bar\eta}(K),
\]
即
\begin{equation}\tag{2.4.5.1}
\operatorname{Var}(\sigma):
\bigl(R\Phi_{\bar\eta}(K)\bigr)_x
\longrightarrow
R\Gamma_{\{x\}}\bigl(R\Psi_{\bar\eta}(K)\bigr).
\end{equation}
原印本在直和项那一句中写的 $R\Phi(F)$, 在此一般表述中应为 $R\Phi(K)$.

\subsubsection*{2.4.6.}
假设 $R\Phi(K)$ 的支集 $\Sigma$ 由有限多个孤立点组成. 则对每个 $\sigma\in I$, $R\Psi_{\bar\eta}(K)$ 的自同态 $\sigma-1$ 是下图中所示的复合:
\begin{equation}\tag{2.4.6.1}
\resizebox{\textwidth}{!}{%
\begin{tikzcd}[ampersand replacement=\&, column sep=huge, row sep=huge]
R\Psi_{\bar\eta}(K)
  \arrow[r] \arrow[d, "\sigma-1"']
  \& R\Phi_{\bar\eta}(K)
     =\displaystyle\bigoplus_{x\in\Sigma}
       i_{x*}\bigl(R\Phi_{\bar\eta}(K)\bigr)_x
     \arrow[d, "{\bigoplus i_{x*}\operatorname{Var}(\sigma)}"'] \\
R\Psi_{\bar\eta}(K)
  \& \displaystyle\bigoplus_{x\in\Sigma}
       i_{x*}R\Gamma_{\{x\}}R\Psi_{\bar\eta}(K)
     \arrow[l]
\end{tikzcd}}
\end{equation}

若 $f$ 固有, 对 (2.4.6.1) 应用 $H^i$, 可见 $\mathbb H^i(X_{\bar\eta},K)$ 的自同态 $\sigma-1$ 为复合
\begin{equation}\tag{2.4.6.2}
\resizebox{\textwidth}{!}{%
\begin{tikzcd}[ampersand replacement=\&, column sep=large, row sep=huge]
\mathbb H^i(X_{\bar\eta},K)
  \arrow[r, "\sim"] \arrow[d, "\sigma-1"']
  \& \mathbb H^i\bigl(X_s,R\Psi_{\bar\eta}(K)\bigr)
     \arrow[r]
  \& \mathbb H^i\bigl(X_s,R\Phi_{\bar\eta}(K)\bigr)
     =\displaystyle\bigoplus_{x\in\Sigma}R^i\Phi_{\bar\eta}(K)_x
     \arrow[d, "{\bigoplus\operatorname{Var}(\sigma)}"'] \\
\mathbb H^i(X_{\bar\eta},K)
  \& \mathbb H^i\bigl(X_s,R\Psi_{\bar\eta}(K)\bigr)
     \arrow[l, "\sim"']
  \& \displaystyle\bigoplus_{x\in\Sigma}
     \mathbb H^i_{\{x\}}\bigl(X_s,R\Psi_{\bar\eta}(K)\bigr)
     \arrow[l]
\end{tikzcd}}
\end{equation}
这描述了 $I$ 在 $\mathbb H^i(X_{\bar\eta},K)$ 上的作用; 其关键组成部分是变差态射 (2.4.5.1), 它们在 $X$ 上是局部的. 还有长正合列
\begin{equation}\tag{2.4.6.3}
\longrightarrow \mathbb H^i(X_s,K)
\longrightarrow \mathbb H^i(X_{\bar\eta},K)
\longrightarrow \bigoplus_{x\in\Sigma}R^i\Phi_{\bar\eta}(K)_x
\longrightarrow\cdots .
\end{equation}

若 $f$ 不必固有, 但 $(X,K)$ 满足 2.1.9, 同样的论证仍然适用. 下图 (2.4.6.5) 定义一个态射
\begin{equation}\tag{2.4.6.4}
\operatorname{Var}(\sigma):
R\Gamma(X_{\bar\eta},K)\longrightarrow R\Gamma_c(X_{\bar\eta},K)
\end{equation}
使 $R\Gamma(X_{\bar\eta},K)$（分别 $R\Gamma_c(X_{\bar\eta},K)$）的自同态 $\sigma-1$ 为复合 $\alpha\operatorname{Var}(\sigma)$（分别 $\operatorname{Var}(\sigma)\alpha$）, 其中
\[
\alpha:R\Gamma_c(X_{\bar\eta},K)\longrightarrow
R\Gamma(X_{\bar\eta},K).
\]
\begin{equation}\tag{2.4.6.5}
\resizebox{\textwidth}{!}{%
\begin{tikzcd}[ampersand replacement=\&, column sep=huge, row sep=huge]
R\Gamma(X_{\bar\eta},K)
  \arrow[r, "\sim"] \arrow[d, "\operatorname{Var}(\sigma)"']
  \& R\Gamma\bigl(X_s,R\Psi_{\bar\eta}(K)\bigr)
     \arrow[r]
  \& \displaystyle\bigoplus_{x\in\Sigma}
     \bigl(R\Phi_{\bar\eta}(K)\bigr)_x
     \arrow[d, "\bigoplus"']
  \& {} \\
R\Gamma_c(X_{\bar\eta},K)
  \arrow[r, "\sim"']
  \& R\Gamma_c\bigl(X_s,R\Psi_{\bar\eta}(K)\bigr)
  \& \displaystyle\bigoplus_{x\in\Sigma}
     R\Gamma_{\{x\}}\bigl(X_s,R\Psi_{\bar\eta}(K)\bigr)
     \arrow[l]
  \& .
\end{tikzcd}}
\end{equation}

在此情形中, 还有长正合列
\begin{equation}\tag{2.4.6.6}
\resizebox{\textwidth}{!}{%
\begin{tikzcd}[ampersand replacement=\&, column sep=large, row sep=large]
\cdots \arrow[r]
  \& \mathbb H^i_c(X_s,K) \arrow[r] \arrow[d]
  \& \mathbb H^i_c(X_{\bar\eta},K) \arrow[r] \arrow[d]
  \& \displaystyle\bigoplus_{x\in\Sigma}
     \bigl(R^i\Phi_{\bar\eta}(K)\bigr)_x
     \arrow[r] \arrow[d, equal]
  \& \cdots \\
\cdots \arrow[r]
  \& \mathbb H^i(X_s,K) \arrow[r]
  \& \mathbb H^i(X_{\bar\eta},K) \arrow[r]
  \& \displaystyle\bigoplus_{x\in\Sigma}
     R^i\Phi_{\bar\eta}(K)_x \arrow[r]
  \& \cdots
\end{tikzcd}}
\end{equation}

在 (2.4.6.2) 之后的讨论中, 原印本所指的 (2.4.4.1) 是变差态射 (2.4.5.1).
